\documentclass[12pt]{article}

\input{../../_assets/preambulo.tex}


\begin{document}

    % 1. Foto de fondo
    % 2. Título
    % 3. Encabezado Izquierdo
    % 4. Color de fondo
    % 5. Coord x del titulo
    % 6. Coord y del titulo
    % 7. Fecha

    
    \input{../../_assets/portada}
    \portadaExamen{etsiitA4.jpg}{LMD\\Examen V}{LMD. Examen V}{MidnightBlue}{-8}{28}{2026}{Roxana Acedo Parra}

    \begin{description}
        \item[Asignatura] Lógica y métodos discretos.
        \item[Curso Académico] 2025-26.
        \item[Grado] Doble Grado en Ingeniería Informática y Matemáticas.
        \item[Grupo] Subgrupo de prácticas.
        \item[Profesor] Francisco Miguel García Olmedo.
        \item[Descripción] Primer parcial.
        \item[Fecha] 25 de marzo de 2026.
        \item[Duración] 60 minutos.
    \end{description}
    
    \newpage

    \begin{ejercicio}[5 puntos] Sea $e$ la función dada por:
        \begin{align}
            e(a,0) &= 1 \\
            e(a,b) &= \begin{cases}
            e\left(a^2, \dfrac{b}{2}\right) & \text{si } b \text{ es par}, \\[0.3cm]
            e\left(a^2, \dfrac{b-1}{2}\right) a & \text{si } b \text{ es impar}.
            \end{cases}
        \end{align}
        Demuestre por inducción que para cualesquiera números naturales $a$ y $b$, $e(a,b) = a^b$.
    \end{ejercicio}
    
    \begin{ejercicio}[5 puntos] Considere el problema de recurrencia:
        $$
        \begin{cases}
        a_0 = 0 \\
        a_1 = -1 \\
        a_n + a_{n-1} + a_{n-2} = n, & \text{para todo } n \ge 2
        \end{cases}
        $$
        y resuélvalo.
    \end{ejercicio}

\end{document}